\section{Related Work}
\label{sec:related-work}

\subsection{Treewidth of Structured Programs}
\label{subsec:rw-treewidth-structured}

% - Thorup (1998): treewidth \leq 6 for goto-free C; linear-time register allocation
% - Bodlaender, Gustedt, Telle (1998): linear-time register allocation for fixed number of registers

\subsection{Treewidth of Java Programs}
\label{subsec:rw-java}

Gustedt, Mæhle, and Telle~\cite{gustedt_treewidth_2002} investigate whether Thorup's result for \texttt{goto}-free C~\cite{thorup_1998} carries over to Java.
Thorup showed that \texttt{goto}-free C programs have CFG treewidth at most~6, using the notion of \emph{flow-affecting constructs} (FACs): language constructs that add non-local control-flow edges and raise treewidth by at most~1 each.

Java has no \texttt{goto}, but its labeled \texttt{break} and \texttt{continue} statements are expressively equivalent to \texttt{goto}.
Consequently, full Java programs can have arbitrarily high treewidth.
Restricting to a label-free subset eliminates this issue: label-free Java has exactly four FACs (\texttt{break}, \texttt{continue}, \texttt{return}, and \texttt{short-circuit evaluation}), giving a treewidth bound of at most~$2 + 4 = 6$.

To complement this theoretical result, the authors implement a tool that computes tree decompositions of Java CFGs and apply it to 9387 methods from the Java standard library.
The average treewidth is~2.7 and the maximum observed is~5, with over 95\% of methods having treewidth at most~3, which is well below the theoretical bound.

This paper is directly relevant to our work: it establishes both a theoretical treewidth bound and empirical baselines for a real-world language, serving as the primary comparison point for our analogous analysis of Rust programs via MIR.

\subsection{Static Analysis of Rust}
\label{subsec:rw-rust-analysis}

% - Prusti: formal verification via MIR
% - Budde (2024): data-flow vulnerability analysis; CFG extraction tooling we may reuse
% - Other MIR-based tools
